% ===== UNIT 07: Cross-domain transfer and dataset bias (2.3) =====
% Status: REWRITTEN  (change to REWRITTEN / REVIEWED when done; see CHECKLIST.md)
% Rules: style rewrite only. Keep every \cite, \ref, \label, number, equation and table value as is. Change the wording, not the content.
\subsection{Cross-domain transfer and dataset bias}
Cross-dataset transfer is a known weak point of ReID systems. Unsupervised domain adaptation (UDA) methods such as SPGAN~\cite{deng2018spgan}, ECN~\cite{zhong2019ecn} and UDAP~\cite{song2020udap} train on a labelled source and adapt using unlabelled target data. They are trained on the large VehicleID source and tested on VeRi-Wild. On the 3{,}000-identity protocol they reach 24.1--34.7\% mAP. The strongest reported direct transfer (no target data at all) reaches 37.7\%. Full viewpoint-aware adaptation reaches 39.7\%~\cite{zheng2020vapc}. Recent work finds that much of the gap between in-domain and cross-domain performance comes from train/test similarity bias inside the ReID benchmarks. Models trained on VeRi-Wild lose ${\approx}15$\,pp mAP when tested on the visually similar VeRi-776. The mAP on vehicle types unseen during training drops from above 90\% to 65--70\%, even for state-of-the-art models~\cite{benmabrouk2026genlimits}. This is the context for the cross-dataset evaluation in Section~\ref{sec:veriwild}. Large drops under zero-shot transfer are normal in this field. They do not come from one particular pipeline.
